\begin{abstract}
This paper studies whether current multimodal and agent-based AI systems can reconstruct machine-usable circuit representations from schematic images. We frame the task as an engineering reconstruction problem rather than a generic vision-language exercise and evaluate three agents, Codex, Claude, and Gemini, within a shared web-based orchestration pipeline. For each input schematic, the pipeline requires dual-modality output generation in the form of Falstad-compatible circuit artefacts and LTSpice schematic artefacts, together with validation-aware reports and rendered preview recovery for comparative inspection. The results show mixed performance. On simpler benchmark circuits, at least partial structural recovery is often possible, and the strongest runs produce recognizable topologies in one or both output modalities. However, performance degrades substantially as circuit complexity increases, with later cases showing sparse Falstad outputs, placeholder LTSpice files, cross-modality inconsistencies, and repeated validation issues. Claude produces the strongest overall artefacts, Codex shows limited intermediate capability, and Gemini most frequently collapses to low-information outputs. The study contributes a practical benchmark workflow for schematic-image reconstruction and provides empirical evidence that current general-purpose agents can sometimes recover useful circuit structure, but do not yet deliver reliable topology-preserving reconstruction without stronger structural modelling and verification.
\end{abstract}
